\section{Introduction}

Asymmetric cell division (ACD) is the process by which cells divide into two daughter cells that differ in fate, size, or both \cite{horvitzMechanismsAsymmetricCell1992, sunchuPrinciplesMechanismsAsymmetric2020, knoblichAsymmetricCellDivision2010, venkei_2018_emergingmechanismsasymmetric, liArtChoreographingAsymmetric2013a}.
During development, ACD allows stem cells to self-renew while producing a population of differentiated cells that will eventually form an organ.
ACD determines the number and identity of cells in developing tissues and organs, thus its regulation is crucial for ensuring consistent developmental outcomes\cite{yang_2021_centrosomeregulationfunction, lechlerSpindlePositioningIts2021, bolkent_2024_cellularmolecularmechanisms}.
In the nervous system, disruption of ACD regulation has been linked to abnormal brain growth and neurodevelopmental disorders, including microcephaly \cite{yang_2021_centrosomeregulationfunction, siskos_2021_moleculargeneticsmicrocephaly, vaidProgenitorBasedCellBiological2022}.

\textit{Drosophila} larval neuroblasts provide a tractable system for studying the regulation of ACD during brain development \cite{gallaud_2017_drosophilamelanogasterneuroblastsa, sunchuPrinciplesMechanismsAsymmetric2020}.
In wild-type flies, a neural stem cell called a neuroblast (NB) divides asymmetrically along a plane perpendicular to its apical/basal axis to produce a larger apical daughter that retains NB identity and a smaller basal daughter called a ganglion mother cell (GMC) (Figure \ref{ch3-fig1} A)\cite{gallaud_2017_drosophilamelanogasterneuroblastsa, cabernard_2009_apicalbasalspindlea, homemDrosophilaNeuroblastsModel2012}.
The GMC then divides once more to generate two terminally differentiated neurons (Figure \ref{ch3-fig1} A)\cite{gallaud_2017_drosophilamelanogasterneuroblastsa, cabernard_2009_apicalbasalspindlea}.
Across successive divisions, this produces a consistent lineage morphology in which the NB remains at one peripheral end of the lineage (Figure \ref{ch3-fig1} c bottom).
This organization suggests that correct neurogenesis depends on both the number of cells produced and the repeated spatial placement of those cells by precisely regulated asymetric division geometry.

Division geometry is not only important in establishing the morphology of the developing neuroblast lineage, but also in determining the identity of the daughter cells \cite{cabernard_2009_apicalbasalspindlea}.
In \textit{mud} mutant flies, spindle orientation is disrupted and approximately 15\% of NB divisions become symmetric, producing two neuroblast daughters side by side instead of a neuroblast and a GMC \cite{sillerNuMArelatedMudProtein2006, cabernard_2009_apicalbasalspindlea}.
This change in the plane of division changes the differential sequestration of cell identity determinant proteins in the daughter cells, leading to differential identity outcomes\cite{cabernard_2009_apicalbasalspindlea}.
Paradoxically, despite having more neuroblasts and therefore greater apparent proliferative potential, \textit{mud} mutant lineages are smaller overall and contain fewer total cells than wild-type lineages \cite{delgado_2025_alteringcellsize}.
Experimental observations further suggest that when multiple neuroblasts are present in a \textit{mud} mutant lineage, they often remain adjacent to one another rather than becoming separated by differentiated progeny.

Many questions about this system are challenging to investigate experimentally.
In particular, several coupled features of asymmetric NB division act simultaneously \invivo: the orientation of the apical/basal axis across generations, stochastic variation in spindle orientation relative to that axis, and the position of the division plane along the axis.
These features are biologically meaningful to separate because they are governed by partially distinct cellular machinery: cortical polarity proteins establish the apical/basal axis, Pins/Mud-dependent machinery aligns the spindle relative to that axis, and polarized actomyosin dynamics together with cleavage-furrow positioning contribute to daughter-size asymmetry and asymmetric furrow (or plane of division) placement \cite{prehodaPolarizationDrosophilaNeuroblasts2009, penkert_2024_drosophilaneuroblastpolarity, sillerNuMArelatedMudProtein2006, cabernard_2010_spindleindependentcleavagefurrow, connell_2011_asymmetriccorticalextension, tsankovaCellPolarityRegulates2017, roubinetSpatiotemporallySeparatedCortical2017a, pham_2019_spatiotemporallycontrolledmyosin}.
Because these features normally change together, it is difficult to determine experimentally which are most responsible for maintaining morphologically consistent wild-type lineages.
It is likewise unclear which regulatory mechanisms are sufficient to prevent \textit{mud} mutant lineages from outgrowing wild type, which mutant division features drive neuroblast adjacency, and what regulatory features could plausibly maintain neuroblast proximity once extra neuroblasts are produced.

Agent-based models (ABMs) are well suited to try to address these questions.
ABMs represent individual cells as interacting entities with explicit rules for growth, division, differentiation, and neighbor-dependent behavior.
Thus they allow us to model and assess how cell behavior may be influenced by internal and external factors.
For a spatial representaation, cellular Potts models (CPMs) are also effective at capturing cell shape, adhesion, packing, and rearrangement, allowing lineage morphology to emerge from local cell-scale interactions\cite{graner_1992_simulationbiologicalcell, glazier_1993_simulationdifferentialadhesion, merks_2005_cellcenteredapproachdevelopmental, hirashima_2017_cellularpottsmodeling}.
This makes CPM-based agent-based models particularly useful for developmental questions in which tissue architecture is an important readout.
These models are powerful tools for studying this system as they allow us to simulate the evolution of the system after imposing perturbations that are difficult or impossible \invivo.

Here, we develop a calibrated CPM-based agent-based model of developing \textit{Drosophila} neuroblast lineages and use it to test how division geometry and growth regulation shape lineage morphology in wild-type and \textit{mud} mutant conditions.
We first calibrate the model to experimental wild-type lineage scale and composition, then ask which separable components of asymmetric division are most important for preserving the sorted wild-type morphology.
We next use the model to evaluate candidate \textit{mud} regulatory mechanisms by testing whether altered division geometry alone explains mutant behavior, whether growth regulation linked to birth-size-dependent division thresholds is required to constrain mutant overgrowth, and whether preferential neuroblast-neuroblast adhesion is necessary to maintain neuroblast proximity.
Our results show that stable intergenerational orientation is more important than moderate spindle-angle variance or reduced division-plane offset for maintaining wild-type lineage sorting, that altered division geometry alone predicts \textit{mud} lineage overgrowth, that volume-based regulation constrains mutant lineages only when division thresholds scale with daughter birth size, and that preferential neuroblast-neuroblast adhesion is not sufficient to keep mutant neuroblasts adjacent.
Together, these findings identify division geometry and size-dependent growth regulation as coupled determinants of reproducible neuroblast lineage architecture.
